\documentclass[11pt,a4paper]{article}

% ------------------------------------------------------------------
%  RUNAWAY : Operation Eclipse - Guide du Joueur & Guide de Deploiement
% ------------------------------------------------------------------
\usepackage[utf8]{inputenc}
\usepackage[T1]{fontenc}
\usepackage[french]{babel}
\usepackage[a4paper,margin=2.3cm]{geometry}
\usepackage{amsmath,amssymb}
\usepackage{xcolor}
\usepackage{booktabs}
\usepackage{array}
\usepackage{longtable}
\usepackage{enumitem}
\usepackage{fancyhdr}
\usepackage{titlesec}
\usepackage[hidelinks,colorlinks=true,linkcolor=eclipsered,urlcolor=eclipseblue]{hyperref}

\definecolor{eclipsered}{RGB}{192,57,43}
\definecolor{eclipseblue}{RGB}{90,120,200}
\definecolor{eclipsedark}{RGB}{30,30,40}
\definecolor{eclipsegreen}{RGB}{60,160,90}
\definecolor{codebg}{RGB}{244,244,248}
\definecolor{tipbg}{RGB}{235,245,238}

\titleformat{\section}{\Large\bfseries\color{eclipsered}}{\thesection.}{0.5em}{}
\titleformat{\subsection}{\large\bfseries\color{eclipsedark}}{\thesubsection}{0.5em}{}

\setlength{\headheight}{14pt}
\pagestyle{fancy}
\fancyhf{}
\lhead{\textcolor{eclipsered}{\textbf{RUNAWAY}} -- Operation Eclipse}
\rhead{Guide v1.0}
\cfoot{\thepage}
\renewcommand{\headrulewidth}{0.4pt}

\newcommand{\code}[1]{\colorbox{codebg}{\texttt{#1}}}
\newcommand{\key}[1]{\fcolorbox{eclipsedark}{white}{\small\texttt{#1}}}

% Encadre "astuce" (capture du contenu via une savebox -> grouping robuste)
\newsavebox{\astucebox}
\newenvironment{astuce}{%
  \begin{lrbox}{\astucebox}%
  \begin{minipage}{\dimexpr0.93\linewidth-2\fboxsep}\color{eclipsedark}%
  \textbf{\textcolor{eclipsegreen}{Astuce}} :\ %
}{%
  \end{minipage}%
  \end{lrbox}%
  \begin{center}\colorbox{tipbg}{\usebox{\astucebox}}\end{center}%
}

\begin{document}

% ------------------------------------------------------------------
\begin{titlepage}
\centering
\vspace*{3cm}
{\Huge\bfseries\color{eclipsered} RUNAWAY}\\[0.4cm]
{\LARGE\color{eclipsedark} Operation Eclipse}\\[1.2cm]
{\Large\bfseries Guide du Joueur \&{} Guide de D\'eploiement}\\[0.6cm]
{\large Version 1.0}\\[2cm]
{\large Jeu d'infiltration top-down --- Godot 4.6.2 / GDScript}\\[0.3cm]
{\large Plateformes : Windows, Linux, Android (Google Play)}\\
\vfill
{\small Document compagnon du \emph{Game Design Document} v4.0}\\
\end{titlepage}

\tableofcontents
\newpage

% ==================================================================
\part*{Partie I --- Guide du Joueur}
\addcontentsline{toc}{section}{Partie I --- Guide du Joueur}

% ------------------------------------------------------------------
\section{Bienvenue, Agent Zero}
Vous incarnez l'\textbf{Agent Zero}, un sp\'ecialiste de l'infiltration envoy\'e
dans le complexe \emph{Eclipse}. Votre objectif global se d\'eroule sur
\textbf{trois niveaux} encha\^in\'es : le Centre de recherche, la For\^et et la
Ville. Chaque niveau demande de r\'ecup\'erer des donn\'ees classifi\'ees puis
de rejoindre la zone d'exfiltration --- id\'ealement \textbf{sans d\'eclencher
l'alarme}.

\medskip
La discr\'etion est r\'ecompens\'ee : un parcours sans alarme et sans dommage
maximise votre score final. Mais le jeu reste jouable \og{}bourrin\fg{} : vous
\^etes arm\'e et pouvez neutraliser les gardes si la situation d\'eg\'en\`ere.

% ------------------------------------------------------------------
\section{Commandes}

\subsection{Clavier \&{} souris (PC)}
\begin{center}
\begin{tabular}{@{}ll@{}}
\toprule
\textbf{Action} & \textbf{Touche} \\
\midrule
Se d\'eplacer        & \key{W} \key{A} \key{S} \key{D} ou \key{Fl\`eches} \\
Sprinter             & \key{Maj} (maintenu) \\
Viser / Tirer        & \key{Clic gauche} (vise la souris) \\
Recharger            & \key{R} \\
Interagir (coffre, terminal) & \key{E} (maintenu) \\
Pause                & \key{\'Echap} \\
\bottomrule
\end{tabular}
\end{center}

\subsection{Tactile (Android)}
\begin{center}
\begin{tabular}{@{}ll@{}}
\toprule
\textbf{Action} & \textbf{Geste} \\
\midrule
Se d\'eplacer    & \textbf{Stick gauche} (pouss\'e \`a fond = sprint) \\
Viser / Tirer    & \textbf{Stick droit} (tir automatique dans la direction vis\'ee) \\
Recharger        & Bouton \key{R} \`a l'\'ecran \\
Interagir        & Bouton \key{E} \`a l'\'ecran (maintenu) \\
Pause            & Bouton \textbf{Retour} du t\'el\'ephone \\
\bottomrule
\end{tabular}
\end{center}

\begin{astuce}
Sur t\'el\'ephone, le \emph{stick gauche} g\`ere \`a la fois la vitesse et le
sprint : un l\'eger d\'eplacement = marche silencieuse, pouss\'e \`a fond =
course bruyante. Dosez votre pouce pour rester discret.
\end{astuce}

% ------------------------------------------------------------------
\section{Lire l'interface (HUD)}
\begin{center}
\begin{tabular}{@{}p{4.2cm}p{10cm}@{}}
\toprule
\textbf{\'El\'ement} & \textbf{Signification} \\
\midrule
Cœurs (haut gauche)        & Points de vie (3 max). Rouge = plein, gris = perdu. \\
Barre bleue (sous les cœurs) & Jauge d'endurance de \textbf{sprint}. Se vide en courant, se recharge \`a l'arr\^et. \\
Munitions (haut droite)    & Balles restantes dans le chargeur (\code{8 / 8}). Affiche \og{}RECHARGEMENT\fg{} pendant le \texttt{reload}. \\
Indicateur d'alarme        & \textcolor{eclipsegreen}{CALME} $\rightarrow$ \textcolor{orange}{SUSPICION} $\rightarrow$ \textcolor{eclipsered}{ALERTE}. \\
Score (bas droite)         & Score cumul\'e de la campagne. \\
Minimap                    & Plan du niveau, position du joueur, objectif et sortie. \\
Banni\`ere de niveau       & \og{}NIVEAU n/3\fg{} affich\'ee au d\'ebut, dispara\^it apr\`es 2\,s. \\
\bottomrule
\end{tabular}
\end{center}

% ------------------------------------------------------------------
\section{Le syst\`eme d'alarme}
L'\'etat d'alerte du complexe \'evolue selon une machine \`a trois \'etats :

\begin{center}
\begin{tabular}{@{}lp{10cm}@{}}
\toprule
\textbf{\'Etat} & \textbf{D\'eclencheur \&{} cons\'equence} \\
\midrule
\textcolor{eclipsegreen}{\textbf{CALME}}     & Situation normale, gardes en patrouille. \\
\textcolor{orange}{\textbf{SUSPICION}} & Un garde vous a entr'aper\c{c}u ou a entendu un bruit. Il enqu\^ete. Si vous vous cachez, le calme revient apr\`es un d\'elai. \\
\textcolor{eclipsered}{\textbf{ALERTE}}      & Vous \^etes rep\'er\'e ou vous avez pirat\'e le coffre. Les gardes vous prennent en chasse et se coordonnent. \\
\bottomrule
\end{tabular}
\end{center}

\begin{astuce}
Une fois le coffre pirat\'e, l'alarme passe \textbf{automatiquement} en ALERTE :
c'est le signal de \textbf{fuir} vers l'exfiltration. Pr\'eparez votre
itin\'eraire de repli \emph{avant} de lancer le piratage.
\end{astuce}

% ------------------------------------------------------------------
\section{L'IA des gardes (ce qu'il faut savoir)}
Les gardes sont \og{}intelligents\fg{} : comprendre leur comportement, c'est
survivre.
\begin{itemize}[leftmargin=1.4em]
  \item \textbf{Vision en c\^one.} Ils voient devant eux dans un c\^one de
  90$^\circ$. La d\'etection est plus fiable au centre du c\^one qu'\`a sa
  p\'eriph\'erie : passer dans le coin de leur champ de vision est moins risqu\'e.
  \item \textbf{Les murs bloquent la vue.} Un \texttt{RayCast} arr\^ete la vision
  derri\`ere murs, caisses et eau. Restez \`a couvert.
  \item \textbf{Le bruit attire.} Sprinter et surtout \textbf{tirer} g\'en\`ere
  du bruit qui alerte les gardes alentour, m\^eme hors de vue.
  \item \textbf{Tir pr\'edictif.} En combat, ils anticipent votre d\'eplacement :
  ne courez pas en ligne droite, zigzaguez.
  \item \textbf{Encerclement.} Plusieurs gardes se r\'epartissent autour de vous
  plut\^ot que de s'agglutiner. \'Evitez de vous laisser entourer : adossez-vous
  \`a un mur ou un couloir.
  \item \textbf{Ils \'evitent le tir ami.} Ils ne tirent pas si un coll\`egue est
  sur la ligne de feu --- vous pouvez exploiter cela en gardant un garde entre
  vous et un autre.
  \item \textbf{Fouille coordonn\'ee.} Quand ils vous perdent, ils fouillent
  votre derni\`ere position connue. Changez de cachette apr\`es avoir bris\'e le
  contact.
\end{itemize}

% ------------------------------------------------------------------
\section{Les trois niveaux}
\begin{center}
\begin{tabular}{@{}clp{7.5cm}@{}}
\toprule
\textbf{\#} & \textbf{Niveau} & \textbf{Ambiance \&{} d\'efi} \\
\midrule
1 & Centre de recherche & Int\'erieur cloisonn\'e, couloirs et bureaux. Apprenez les bases de la furtivit\'e. \\
2 & La For\^et          & Espaces ouverts, v\'eg\'etation et points d'eau. Plus de gardes mobiles. \\
3 & La Ville            & Environnement urbain dense, exfiltration par h\'eliport. Mont\'ee finale en tension. \\
\bottomrule
\end{tabular}
\end{center}

\noindent Boucle de jeu de chaque niveau :
\begin{enumerate}[leftmargin=1.6em]
  \item Infiltrez-vous en \'evitant les c\^ones de vision.
  \item Atteignez le coffre / terminal et maintenez \key{E} pour pirater.
  \item L'alarme se d\'eclenche : foncez vers la zone d'exfiltration.
  \item La sortie ne s'ouvre que si l'objectif est rempli (\code{needs\_objective}).
  \item L'exfiltration r\'eussie encha\^ine automatiquement sur le niveau suivant ;
  votre score, vos vies et votre temps sont \textbf{cumul\'es}.
\end{enumerate}

% ------------------------------------------------------------------
\section{Maximiser son score}
Le score final r\'ecompense l'efficacit\'e et la discr\'etion :
\begin{itemize}[leftmargin=1.4em]
  \item \textbf{Discr\'etion totale} (aucune alarme d\'eclench\'ee sur toute la
  campagne) = gros bonus.
  \item \textbf{Rapidit\'e} : un temps de mission court rapporte plus.
  \item \textbf{Int\'egrit\'e} : terminer avec un maximum de vies.
  \item Neutraliser des gardes rapporte des points, mais l'approche
  \og{}fant\^ome\fg{} (sans tir) reste la plus payante.
\end{itemize}

\begin{astuce}
Pour un \emph{run} parfait : marchez (jamais de sprint inutile), tirez le moins
possible, et m\'emorisez les patrouilles avant de bouger. Le silence vaut plus
de points qu'un tableau de chasse.
\end{astuce}

\newpage
% ==================================================================
\part*{Partie II --- Guide de Build \&{} D\'eploiement}
\addcontentsline{toc}{section}{Partie II --- Guide de Build \&{} D\'eploiement}

% ------------------------------------------------------------------
\section{Pr\'erequis}
\begin{center}
\begin{tabular}{@{}lp{9.5cm}@{}}
\toprule
\textbf{Outil} & \textbf{Usage} \\
\midrule
Godot \textbf{4.6.2} (.NET non requis) & Moteur et exporteur. \\
\emph{Android Build Template} & Mod\`ele Gradle export\'e depuis Godot. \\
JDK 17 + Android SDK          & Compilation Gradle de l'APK/AAB. \\
\code{keytool} (fourni avec le JDK) & G\'en\'eration du keystore de signature. \\
PowerShell                    & Ex\'ecution des scripts utilitaires (\code{tools/}). \\
\bottomrule
\end{tabular}
\end{center}

% ------------------------------------------------------------------
\section{\'Etape 1 --- Pr\'eparer le projet}
\begin{enumerate}[leftmargin=1.6em]
  \item Ouvrez le projet dans Godot 4.6.2.
  \item Menu \textbf{Projet $\rightarrow$ Installer le mod\`ele de build Android}
  (\emph{Install Android Build Template}). Cela cr\'ee le dossier \code{android/build/}.
  \item Dans \textbf{\'Editeur $\rightarrow$ Param\`etres de l'\'editeur $\rightarrow$
  Export $\rightarrow$ Android}, renseignez les chemins du \textbf{JDK}, du
  \textbf{SDK Android} et de \code{adb}/\code{keytool}.
\end{enumerate}

\noindent Les ic\^ones de lanceur sont d\'ej\`a g\'en\'er\'ees dans
\code{mobile\_icons/} (cible rouge sur fond sombre). Pour les r\'eg\'en\'erer :
\begin{quote}\code{python tools/gen\_icons.py}\end{quote}

% ------------------------------------------------------------------
\section{\'Etape 2 --- Cr\'eer le keystore de signature}
Google Play exige un APK/AAB \textbf{sign\'e}. Un script automatise la cr\'eation
du keystore de \emph{release} et l'enregistrement des variables d'environnement
attendues par Godot :
\begin{quote}
\code{pwsh -File tools/gen\_keystore.ps1 -Password "VOTRE\_MOT\_DE\_PASSE"}
\end{quote}
Ce script :
\begin{itemize}[leftmargin=1.4em]
  \item cr\'ee \code{\%USERPROFILE\%\textbackslash runaway-release.keystore} (alias \code{eclipse}) ;
  \item d\'efinit les variables utilisateur \code{GODOT\_ANDROID\_KEYSTORE\_RELEASE\_PATH},
  \code{\_USER} et \code{\_PASSWORD}.
\end{itemize}

\begin{astuce}
\textbf{Conservez ce keystore et son mot de passe en lieu s\^ur.} Google Play
associe d\'efinitivement votre application \`a cette cl\'e : la perdre vous
emp\^eche de publier toute mise \`a jour. Le fichier est volontairement exclu du
d\'ep\^ot via \code{.gitignore}.
\end{astuce}

% ------------------------------------------------------------------
\section{\'Etape 3 --- Exporter}
Deux pr\'esets sont d\'efinis dans \code{export\_presets.cfg} :
\begin{center}
\begin{tabular}{@{}llp{6cm}@{}}
\toprule
\textbf{Pr\'eset} & \textbf{Format} & \textbf{Usage} \\
\midrule
Android APK            & \code{.apk} & Tests sur appareil / \emph{sideload}. \\
Android AAB (Play Store) & \code{.aab} & Publication sur Google Play Console. \\
\bottomrule
\end{tabular}
\end{center}

\noindent Caract\'eristiques communes : build Gradle activ\'e, architectures
\code{arm64-v8a} + \code{armeabi-v7a}, \code{versionCode=1}, \code{versionName="1.0"},
identifiant \code{com.eclipse.runaway}, mode immersif, signature \emph{release}
activ\'ee, et \code{exclude\_filter} retirant \code{tools/}, \code{*.py},
\code{mockups/} et \code{*.md} du paquet.

\medskip
\noindent\textbf{Depuis l'\'editeur :} Projet $\rightarrow$ Exporter $\rightarrow$
choisir le pr\'eset $\rightarrow$ \emph{Exporter le projet}.

\medskip
\noindent\textbf{En ligne de commande (headless) :}
\begin{quote}
\code{godot --headless --export-release "Android AAB (Play Store)" build/runaway.aab}\\[2pt]
\code{godot --headless --export-release "Android APK" build/runaway.apk}
\end{quote}

% ------------------------------------------------------------------
\section{\'Etape 4 --- Publier sur Google Play}
\begin{enumerate}[leftmargin=1.6em]
  \item Cr\'eez un compte d\'eveloppeur sur la \textbf{Google Play Console}.
  \item Cr\'eez une nouvelle application, renseignez nom, langue et cat\'egorie
  (\og{}Jeu / Action\fg{}).
  \item Compl\'etez la \textbf{fiche Play Store} : description, captures d'\'ecran
  (paysage), ic\^one 512$\times$512, banni\`ere de fonctionnalit\'e.
  \item Remplissez les questionnaires obligatoires : \textbf{classification du
  contenu} (PEGI/IARC), \textbf{politique de confidentialit\'e}, \textbf{s\'ecurit\'e
  des donn\'ees}, et \textbf{public cible}.
  \item Cr\'eez une \textbf{version} sur une piste (interne, ferm\'ee, ouverte ou
  production) et t\'el\'eversez le fichier \code{build/runaway.aab}.
  \item Soumettez pour \textbf{examen}. Apr\`es validation, publiez.
\end{enumerate}

% ------------------------------------------------------------------
\section{Conformit\'e Play Store (checklist)}
\begin{center}
\begin{tabular}{@{}p{8.5cm}c@{}}
\toprule
\textbf{Crit\`ere} & \textbf{\'Etat} \\
\midrule
Format AAB pour publication                       & \textcolor{eclipsegreen}{OK} \\
Signature \emph{release} (keystore d\'edi\'e)     & \textcolor{eclipsegreen}{OK} \\
64 bits (\code{arm64-v8a}) inclus                 & \textcolor{eclipsegreen}{OK} \\
Permissions minimales (aucune permission sensible) & \textcolor{eclipsegreen}{OK} \\
Ic\^ones adaptatives (fg/bg 432\,px)              & \textcolor{eclipsegreen}{OK} \\
Orientation paysage verrouill\'ee                 & \textcolor{eclipsegreen}{OK} \\
Bouton \textbf{Retour} g\'er\'e (pause, pas de sortie brutale) & \textcolor{eclipsegreen}{OK} \\
Interface adapt\'ee au tactile (sticks, boutons)  & \textcolor{eclipsegreen}{OK} \\
\code{versionCode} incr\'ementable \`a chaque MAJ & \`a g\'erer \\
\bottomrule
\end{tabular}
\end{center}

\begin{astuce}
\`A chaque nouvelle publication, \textbf{incr\'ementez \code{versionCode}}
(1, 2, 3, \dots) dans le pr\'eset d'export : Google Play refuse deux
t\'el\'eversements avec le m\^eme code de version.
\end{astuce}

% ------------------------------------------------------------------
\section{D\'epannage}
\begin{center}
\begin{tabular}{@{}p{6.5cm}p{8cm}@{}}
\toprule
\textbf{Sympt\^ome} & \textbf{Solution} \\
\midrule
\og{}Android build template not installed\fg{} & Refaire l'\'etape 1 (installer le mod\`ele). \\
\'Echec de signature                            & V\'erifier les variables d'env. du keystore (\code{gen\_keystore.ps1}). \\
\og{}1 resource still in use at exit\fg{} (headless) & Avertissement b\'enin (musique en boucle de \code{AudioManager}), sans cons\'equence. \\
APK rejet\'e (32 bits seul)                     & V\'erifier que \code{arm64-v8a} est coch\'e dans le pr\'eset. \\
\bottomrule
\end{tabular}
\end{center}

\vfill
\begin{center}\small
\textcolor{eclipsedark}{Fin du guide --- bonne infiltration, Agent Zero.}
\end{center}

\end{document}
